\section{Zaključak}

U radu je razvijen i primijenjen funkcionalan sustav za mjerenje trodimenzijskih
pomaka i deformacija površine objekta pomoću stereo korelacije digitalne slike.
Primijenjen je pristup koji spaja globalnu 2D korelaciju konačnih elemenata po svakoj
kameri sa stereo triangulacijom čvorova mreže, čime je uz minimalnu količinu novoga
koda ostvarena potpuna 3D rekonstrukcija.

Sustav je uspješno primijenjen na stvarni vlačni eksperiment 3D tiskane polimerne
epruvete s 300 deformacijskih stanja. Svi su rezultati, uz pomoć mjerila slike od
$0{,}0502$\,mm/px i iz njega izvedenog mjerila prostora, izraženi u fizikalnim
jedinicama. Rekonstruirana su glatka polja trodimenzijskih pomaka, jasno
obuhvaćajući i ravninske
i izvanravninske komponente. U posljednjem stanju najveći ukupni pomak iznosi
$0{,}97$\,mm, a najveća izvanravninska komponenta $0{,}87$\,mm, dakle izvanravninski
pomak gotovo je jednak ravninskom. Pripadna polja površinskih deformacija pokazuju
medijalnu glavnu deformaciju od približno 2\,\%.

Mjerna nesigurnost sustava određena je izravno iz 20 slika neopterećenog uzorka.
Standardna nesigurnost pomaka iznosi $2{,}57\cdot10^{-3}$\,mm ($0{,}05$\,px) za
kameru 1, odnosno $3{,}40\cdot10^{-3}$\,mm ($0{,}07$\,px) za kameru 2, uz
nesigurnost deformacija reda veličine $10^{-3}$. Analiza ovisnosti nesigurnosti
o veličini elementa potvrdila je da je odabrana mreža nominalne veličine $10$\,px
povoljan kompromis razlučivosti i šuma, a ispravnost implementacije potvrđena je
rekonstrukcijom referentnog oblika do
strojne preciznosti i subpikselnom reprojekcijskom pogreškom ($\approx 0{,}15$\,px).
Lokalizirani gubitak korelacije na tri rubna čvora jedne kamere automatski je
detektiran usporedbom sa susjednim čvorovima i saniran robusnim medijalnim filtrom,
bez utjecaja na ostatak mjernog polja i bez izmjena mjernih algoritama.

Time su ispunjeni svi ciljevi zadatka. Postavljen je stereovizijski model
(kalibriran ChArUco pločom), provedena
korelacija između referentnih i deformiranih slika, izračunati 3D pomaci i
deformacije u stvarnim jedinicama, određena mjerna nesigurnost te vizualizirano
deformirano stanje. Mogući smjerovi daljnjeg rada
uključuju izglađivanje polja deformacija ekstrapolacijom u čvorove, profinjenje mreže
uz rub mjerne zone te inkrementalnu strategiju korelacije (ažuriranje referentne
slike) kojom bi se analiza proširila i na preostala deformacijska stanja do loma.
